\section{Introduction}\label{sec:intro}

The rapid evolution of Large Language Model (LLM) coding agents from theoretical concepts to practical products between 2023 and 2025 has driven significant progress in automated software engineering. As demonstrated by standard benchmarks like SWE-bench~\cite{jimenez2024swebench}, resolution rates on real-world software engineering tasks have surged from a mere 12\% to 43\%~\cite{yang2024sweagent}. However, recent research~\cite{wang2026patchdiff} highlights lingering reliability issues, particularly regarding behavioral deviations in patches that are otherwise marked as ``resolved.'' These deviations pose catastrophic risks in high-stakes domains such as decentralized finance (DeFi), where the Total Value Locked (TVL) exceeds \$50 billion. The frequent and devastating financial losses caused by vulnerabilities in traditional Solidity contracts---such as the infamous DAO hack (\$60 million) and the Poly Network exploit (\$611 million)---have proven that manual auditing is both inefficient and prohibitively expensive.

To mitigate these risks, the Move programming language was introduced~\cite{blackshear2019move}. With its resource-oriented design, Move provides a superior security baseline and is inherently more friendly to formal verification. Despite these architectural advantages, statistical data from platforms like MoveScan~\cite{song2024movescan} indicates that defect management remains a persistent challenge: among 37,302 deployed contracts and 97,028 recorded defects, the Move Prover automatically catches only 6.02\%. A poignant example is the Thala Labs incident in November 2024, where a seemingly simple 2-line code modification led to a \$25.5 million exploit~\cite{halborn2024thala}. This underscores a critical reality: human intuition is insufficient, and mathematically guaranteed correctness through formal verification is indispensable. However, bridging the gap between high-level specifications and verified code remains incredibly difficult. While existing LLM-assisted tools for Move (such as MSG~\cite{fu2025msg}) focus primarily on generating specifications from existing code (Code$\rightarrow$Spec), our research tackles the inverse and significantly more challenging direction: generating verified contract bodies directly from specifications (Spec$\rightarrow$Code).

Program synthesis presents a promising solution by treating code generation as a search problem constrained by verification. Traditional hybrid synthesis tools attempt to combine deductive and inductive methods; for instance, on 27 Solidity contracts, SmartSpec~\cite{smartspec} achieves an 85\% semantic match. However, these legacy synthesizers are heavily coupled with traditional imperative languages like Solidity. When confronted with resource-oriented languages like Move---which enforce strict ownership semantics and intricate prover constraints---they completely fall short, being further hampered by heavy specification burdens and inefficient feedback loops. Conversely, while LLMs excel at translating natural language into code, they struggle significantly with domain-specific rigors, such as Aptos Move verification. When a single LLM attempts to both generate code and parse obscure, machine-level verification errors to fix its own mistakes, it falls prey to two documented failure modes. First, Huang \emph{et al.}~\cite{huang2024llmselfcorrect} show that LLMs cannot self-correct reasoning without an external oracle: a model reviewing its own output lacks the ground-truth signal needed to distinguish genuine fixes from superficial ones. Second, Wei \emph{et al.}~\cite{wei2024sycophancy,wei2023syntheticsycophancy} demonstrate that LLMs exhibit \emph{sycophancy}---over-conditioning on their own prior outputs and user context, which leads them to ``hallucinate'' logic that superficially satisfies feedback without addressing the root cause. In our setting, the model's previously generated (flawed) body becomes the very context that biases subsequent repair attempts, causing the feedback loop to diverge rather than converge.

To understand why raw prover feedback is insufficient, we examine a concrete synthesis task from the Aptos \texttt{stake} module: the \texttt{update\_performance\_statistics} function. This case exposes two critical gaps: verifier pass alone can coincide with semantic omission under partial specifications, and generic feedback loops fail to converge because prover errors do not directly map to the idiom-level fixes the model needs. A detailed walkthrough of this motivating example is given in \S\ref{sec:motivating-example}.

Driven by these insights, we propose \textbf{MoVES} (\textbf{Mo}ve \textbf{V}erified \textbf{E}xecution \textbf{S}ynthesis), a novel automated synthesis system built upon a formal verification feedback loop that systematizes and automates this structured diagnostic process. MoVES implements a dual-role architecture that explicitly divides the task into an independent \textit{Code Generation Role} and an \textit{Error Diagnosis Role}. To validate this approach, we evaluate three state-of-the-art LLMs on the same synthesis tasks. Under zero-shot settings and generic feedback, Kimi fails entirely to synthesize valid code. Claude Opus 4.7, acting as the generator in our MoVES framework, successfully passes the verification after one round of structured diagnosis (72.7 seconds). Interestingly, while GPT 5.5 appears to pass the test in a zero-shot setting, a closer inspection reveals that it entirely omits the \texttt{failed\_proposer\_indices} logic. This exposes a critical systemic flaw in the commonly used Pass@1 metric: verifier success alone does not guarantee semantic completeness, especially given that real-world specifications are often partial and focus strictly on critical security invariants rather than full functional equivalence.

% TODO: Write contributions and paper organisation after the rest of the paper is complete.
